\documentclass[11pt,a4paper]{article}

% ── Packages ─────────────────────────────────────────────────────────────────
\usepackage[utf8]{inputenc}
\usepackage[T1]{fontenc}
\usepackage[margin=1in]{geometry}
\usepackage{amsmath,amssymb}
\usepackage{booktabs}
\usepackage{array}
\usepackage{listings}
\usepackage{xcolor}
\usepackage{parskip}
\usepackage{hyperref}
\usepackage{caption}
\usepackage{float}
\usepackage{microtype}
\usepackage{fancyhdr}
\usepackage{titlesec}
\usepackage{enumitem}

% ── Page style ───────────────────────────────────────────────────────────────
\pagestyle{fancy}
\fancyhf{}
\rhead{Assignment 4 -- Graph Optimization}
\lhead{\leftmark}
\cfoot{\thepage}

% ── Section formatting ────────────────────────────────────────────────────────
\titleformat{\section}[block]{\large\bfseries}{}{0pt}{}[\titlerule]
\titleformat{\subsection}[block]{\normalsize\bfseries}{}{0pt}{}

% ── Code listing style ────────────────────────────────────────────────────────
\definecolor{codebg}{rgb}{0.96,0.96,0.96}
\definecolor{codegreen}{rgb}{0.10,0.50,0.10}
\definecolor{codegray}{rgb}{0.50,0.50,0.50}
\definecolor{codeblue}{rgb}{0.10,0.10,0.75}
\definecolor{codepurple}{rgb}{0.55,0.00,0.55}

\lstdefinestyle{python}{
  language=Python,
  backgroundcolor=\color{codebg},
  basicstyle=\ttfamily\footnotesize,
  keywordstyle=\color{codeblue}\bfseries,
  commentstyle=\color{codegreen}\itshape,
  stringstyle=\color{codepurple},
  numberstyle=\tiny\color{codegray},
  numbers=left,
  numbersep=6pt,
  showstringspaces=false,
  breaklines=true,
  frame=single,
  framesep=4pt,
  tabsize=4,
  captionpos=b,
}

\lstdefinestyle{output}{
  backgroundcolor=\color{codebg},
  basicstyle=\ttfamily\footnotesize,
  breaklines=true,
  frame=single,
  framesep=4pt,
  numbers=none,
  captionpos=b,
}

\lstset{style=python}

% ── Document ─────────────────────────────────────────────────────────────────
\begin{document}

\begin{titlepage}
  \centering
  \vspace*{2cm}
  {\Huge\bfseries Assignment 4\par}
  \vspace{0.5cm}
  {\LARGE Graph Optimization, Greedy Methods,\\[4pt] and Dynamic Programming\par}
  \vspace{1.5cm}
  {\large Algorithms --- Spring 2026\par}
  \vspace{0.5cm}
  {\large\today\par}
  \vfill
  \textit{Theme: Resilient logistics and infrastructure planning\\
  for a multi-campus university network.}
\end{titlepage}

\tableofcontents
\newpage

% ═══════════════════════════════════════════════════════════════════════════════
\section{Part A --- Theory and Algorithmic Reasoning}
% ═══════════════════════════════════════════════════════════════════════════════

\subsection{A1 --- Graph Representation}

\paragraph{Case 1: Sparse graph with $V = 100{,}000$ vertices and $E = 300{,}000$ edges.}

An \textbf{adjacency list} is the appropriate choice.

\begin{itemize}[noitemsep]
  \item \textbf{Memory.} An adjacency list consumes $O(V + E) = O(400{,}000)$ space. An adjacency matrix would require $O(V^2) = O(10^{10})$ entries -- roughly 75\,GB for 64-bit integers. The list representation is orders of magnitude more compact.
  \item \textbf{Edge lookup.} $O(\deg(u))$ per query in a list vs.\ $O(1)$ in a matrix. For sparse graphs $\deg(u) \ll V$, so the list is still fast enough in practice.
  \item \textbf{Neighbour traversal.} $O(\deg(u))$ for both representations. In the sparse case, the list visits only the actual neighbours rather than scanning all $V$ columns.
  \item \textbf{Overall suitability.} All classic sparse-graph algorithms (BFS, DFS, Dijkstra, Bellman-Ford, Prim) are designed around adjacency lists. The matrix is simply infeasible at this scale.
\end{itemize}

\paragraph{Case 2: Dense graph with $V = 4{,}000$ vertices.}

An \textbf{adjacency matrix} is the appropriate choice.

\begin{itemize}[noitemsep]
  \item \textbf{Memory.} $O(V^2) = O(16 \times 10^6)$ entries -- about 128\,MB for 64-bit values, which is manageable.  A dense graph has $E \approx V^2$, so the adjacency list uses roughly the same memory without the $O(1)$-lookup benefit.
  \item \textbf{Edge lookup.} $O(1)$ in the matrix: check \texttt{A[u][v]} directly. Useful for algorithms like Floyd-Warshall.
  \item \textbf{Neighbour traversal.} $O(V)$ in both representations when the graph is dense (nearly every column is non-zero anyway).
  \item \textbf{Overall suitability.} When $E \sim V^2$, the matrix's constant-time lookup and cache locality make it preferable.
\end{itemize}

\begin{table}[H]
\centering
\caption{Adjacency List vs.\ Adjacency Matrix comparison.}
\begin{tabular}{lll}
\toprule
Criterion & Adjacency List & Adjacency Matrix \\
\midrule
Space & $O(V+E)$ & $O(V^2)$ \\
Edge lookup & $O(\deg(u))$ & $O(1)$ \\
Neighbour traversal & $O(\deg(u))$ & $O(V)$ \\
Best when & Sparse graph ($E \ll V^2$) & Dense graph ($E \sim V^2$) \\
\bottomrule
\end{tabular}
\end{table}

% ─────────────────────────────────────────────────────────────────────────────
\subsection{A2 --- BFS vs.\ Weighted Shortest Path}

\paragraph{When BFS returns a correct shortest path.}
BFS computes shortest paths \emph{in terms of hop count} (number of edges traversed). It is therefore correct whenever all edge weights are equal -- most importantly, on \emph{unweighted} graphs where every edge has implicit weight~1.

\paragraph{When BFS does not return a correct shortest path.}
BFS is incorrect on \emph{weighted} graphs where edges carry different costs. Because BFS expands vertices in order of hop count rather than path weight, it can ``commit'' to a path with fewer hops that has a higher total cost.

\paragraph{Small example.}
Consider the graph:
\[
S \xrightarrow{1} A,\quad S \xrightarrow{10} B,\quad A \xrightarrow{1} B.
\]
BFS from $S$: both $A$ and $B$ are reachable in 1 hop directly from $S$.  BFS reports $\mathrm{dist}(S,B) = 1$ (direct edge).  But the weighted shortest path is $S \to A \to B$ with cost $1+1=2$, which is less than the direct edge weight of 10.  Here BFS gives the wrong \emph{weighted} cost while being correct about hop count.

Conversely, if we ask for the minimum-weight path in this graph, BFS (on the unweighted version) reports $S \to B$ (1 hop) but the weight-optimal path is $S \to A \to B$ (cost 2 vs.\ 10). BFS cannot detect this because it ignores edge weights.

% ─────────────────────────────────────────────────────────────────────────────
\subsection{A3 --- Dijkstra Failure Case with Negative Edges}

\paragraph{Constructed graph.}
\[
S \xrightarrow{3} A,\quad S \xrightarrow{4} B,\quad B \xrightarrow{-5} A.
\]

\paragraph{Dijkstra's behaviour (incorrect).}
\begin{enumerate}[noitemsep]
  \item Initialise: $d[S]=0$, $d[A]=\infty$, $d[B]=\infty$.
  \item Pop $S$: relax $S \to A$ ($d[A]=3$) and $S \to B$ ($d[B]=4$). Heap: $\{(3,A),(4,B)\}$.
  \item Pop $A$ (smallest key $= 3$): $A$ is \textbf{finalised} with distance 3. $A$ has no outgoing edges.
  \item Pop $B$: relax $B \to A$: $d[A] = \min(3,\, 4+(-5)) = \min(3,-1) = -1$. However, $A$ was already removed from the priority queue and marked as finalised, so this update is never propagated.
  \item Dijkstra reports $d[A] = 3$.  The correct answer is $-1$ (via $S \to B \to A$).
\end{enumerate}

\paragraph{Why Dijkstra fails.}
Dijkstra's greedy invariant assumes that once a vertex is extracted from the min-heap its tentative distance is optimal. Negative edges violate this: a later, seemingly longer path can ``shortcut'' through a negative edge and arrive cheaper.

\paragraph{Why Bellman-Ford works.}
Bellman-Ford relaxes \emph{all} edges in every pass, with no finalisation step. After $|V|-1$ passes it is guaranteed to have propagated every shortest path (assuming no negative cycle). On the same graph:
\begin{itemize}[noitemsep]
  \item Pass 1: relax $S\to A$ ($d[A]=3$), $S\to B$ ($d[B]=4$), $B\to A$ ($d[A]=-1$).
  \item Pass 2: no further improvement.
\end{itemize}
Bellman-Ford correctly reports $d[A]=-1$.

% ─────────────────────────────────────────────────────────────────────────────
\subsection{A4 --- Greedy vs.\ Dynamic Programming}

\paragraph{Problem where greedy is optimal: Activity Selection (maximise job count).}
Sort jobs by finish time and greedily accept a job whenever it does not overlap with the last accepted job. This is provably optimal for \emph{maximising the number} of non-overlapping jobs. The exchange argument shows that replacing any greedily chosen job with another can only worsen or maintain the count (see Part~D1 for the implementation).

\paragraph{Superficially similar problem where greedy fails: Weighted Interval Scheduling (maximise total reward).}

Consider the jobs:
\[
A = [1,2],\; \text{reward} = 1 \qquad B = [1,10],\; \text{reward} = 100.
\]
Both start at time 1 and are mutually incompatible. Sorted by finish time: $A$ (finishes at 2) before $B$ (finishes at 10).

\begin{itemize}[noitemsep]
  \item \textbf{Greedy (earliest finish):} selects $A$ (last finish $= 2$), then checks $B$: $\mathrm{start}(B) = 1 < 2$ -- skip. Result: $\{A\}$, reward~$= 1$.
  \item \textbf{DP optimal:} select only $B$, reward $= 100$.
\end{itemize}

The greedy strategy picks the job that finishes earliest, but that job has negligible reward compared to a longer, higher-value job. Only DP considers all combinations and finds the globally optimal subset.

% ═══════════════════════════════════════════════════════════════════════════════
\section{Part B --- Graph Algorithms}
% ═══════════════════════════════════════════════════════════════════════════════

All Part~B algorithms operate on \texttt{data/roads.csv}: a directed weighted graph with vertices $\{A,B,C,D\}$ and edges:
\[
A\!\to\!B\ (4),\quad A\!\to\!C\ (2),\quad C\!\to\!B\ (1),\quad B\!\to\!D\ (5),\quad C\!\to\!D\ (8).
\]

% ─────────────────────────────────────────────────────────────────────────────
\subsection{B1 --- BFS: Unweighted Reachability and Hop Distance}

\paragraph{Approach.}  Edge weights are ignored; the graph is treated as unweighted. BFS from source $A$ computes the minimum number of hops to every reachable vertex. Parent pointers enable path reconstruction in $O(|path|)$.

\begin{lstlisting}[style=python, caption={Core BFS logic (src/b1\_bfs.py).}]
def bfs(graph, vertices, source):
    dist   = {v: -1   for v in vertices}
    parent = {v: None for v in vertices}
    dist[source] = 0
    queue = deque([source])

    while queue:
        u = queue.popleft()
        for v in graph.get(u, []):
            if dist[v] == -1:          # not yet visited
                dist[v]   = dist[u] + 1
                parent[v] = u
                queue.append(v)

    return dist, parent


def reconstruct_path(parent, source, target):
    path = []
    cur  = target
    while cur is not None:
        path.append(cur)
        cur = parent[cur]
    path.reverse()
    return path if (path and path[0] == source) else None
\end{lstlisting}

\paragraph{Output.}
\begin{lstlisting}[style=output]
BFS from source: A
Target     Hops       Path
---------------------------------------------
A          0          A
B          1          A -> B
C          1          A -> C
D          2          A -> B -> D
\end{lstlisting}

Both $B$ and $C$ are one hop from $A$. $D$ requires two hops (via $B$). All four vertices are reachable.

% ─────────────────────────────────────────────────────────────────────────────
\subsection{B2 --- DFS: Structure Analysis}

\paragraph{Approach.}  DFS assigns discovery and finish timestamps to each vertex and classifies every edge according to its relationship in the DFS forest.

\begin{lstlisting}[style=python, caption={DFS visit routine (src/b2\_dfs.py).}]
def dfs_visit(u):
    color[u] = 'gray';  time[0] += 1;  disc[u] = time[0]

    for v in graph.get(u, []):
        if color[v] == 'white':
            parent[v] = u
            edge_classes.append((u, v, 'TREE'))
            dfs_visit(v)
        elif color[v] == 'gray':
            edge_classes.append((u, v, 'BACK'))
        else:                         # black
            if disc[u] < disc[v]:
                edge_classes.append((u, v, 'FORWARD'))
            else:
                edge_classes.append((u, v, 'CROSS'))

    color[u] = 'black';  time[0] += 1;  fin[u] = time[0]
\end{lstlisting}

\paragraph{Output.}
\begin{lstlisting}[style=output]
DFS Structure Analysis
================================================
Vertex     disc     fin      parent
------------------------------------------------
A          1        8        None
B          2        5        A
C          6        7        A
D          3        4        B

Edge Classification:
-----------------------------------
  A -> B  :  TREE
  B -> D  :  TREE
  A -> C  :  TREE
  C -> B  :  CROSS
  C -> D  :  CROSS

Discussion:
  No back edges: graph is a DAG (no directed cycles).
  Cross edges [('C','B'),('C','D')]: reach already-finished subtrees,
  confirming the graph has multiple paths without creating cycles.
  All 4 vertices discovered from one DFS root -> weakly connected.
\end{lstlisting}

\paragraph{Discussion.}
\begin{itemize}[noitemsep]
  \item \textbf{Cycles.} No back edges are found, so the graph is a directed acyclic graph (DAG). This is consistent with the road network having no circular routes.
  \item \textbf{Connectivity.} DFS from $A$ (the alphabetically first vertex) reaches all four vertices without needing to restart -- the graph is weakly connected.
  \item \textbf{Traversal order.} $A$ is the root; $B$ and its subtree ($D$) are explored first ($A\to B\to D$); $C$ is explored last, contributing two cross edges to already-finished vertices $B$ and $D$. Cross edges confirm alternative paths between vertices that do not create cycles.
\end{itemize}

% ─────────────────────────────────────────────────────────────────────────────
\subsection{B3 --- Dijkstra and Bellman-Ford}

\begin{lstlisting}[style=python, caption={Dijkstra's algorithm (src/b3\_shortest\_paths.py).}]
def dijkstra(graph, vertices, source):
    dist   = {v: float('inf') for v in vertices}
    parent = {v: None         for v in vertices}
    dist[source] = 0
    pq = [(0, source)]

    while pq:
        d, u = heapq.heappop(pq)
        if d > dist[u]:        # stale heap entry
            continue
        for v, w in graph.get(u, []):
            new_dist = dist[u] + w
            if new_dist < dist[v]:
                dist[v]   = new_dist
                parent[v] = u
                heapq.heappush(pq, (new_dist, v))

    return dist, parent
\end{lstlisting}

\begin{lstlisting}[style=python, caption={Bellman-Ford algorithm (src/b3\_shortest\_paths.py).}]
def bellman_ford(edges, vertices, source):
    dist   = {v: float('inf') for v in vertices}
    parent = {v: None         for v in vertices}
    dist[source] = 0

    n = len(vertices)
    for _ in range(n - 1):
        updated = False
        for u, v, w in edges:
            if dist[u] != float('inf') and dist[u] + w < dist[v]:
                dist[v]   = dist[u] + w
                parent[v] = u
                updated   = True
        if not updated:        # converged early
            break

    # One extra pass: any improvement means a negative cycle exists.
    has_neg_cycle = any(
        dist[u] != float('inf') and dist[u] + w < dist[v]
        for u, v, w in edges
    )
    return dist, parent, has_neg_cycle
\end{lstlisting}

\paragraph{Output.}
\begin{lstlisting}[style=output]
Dijkstra – shortest paths from 'A':
Target     Distance     Path
--------------------------------------------------
A          0            A
B          3            A -> C -> B
C          2            A -> C
D          8            A -> C -> B -> D

Bellman-Ford – shortest paths from 'A':
Target     Distance     Path
--------------------------------------------------
A          0            A
B          3            A -> C -> B
C          2            A -> C
D          8            A -> C -> B -> D

Negative cycle detected: False
\end{lstlisting}

\paragraph{Discussion.}
Both algorithms agree because \texttt{roads.csv} has no negative weights.

\begin{itemize}[noitemsep]
  \item \textbf{Dijkstra} runs in $O((V+E)\log V)$ with a binary heap. Its greedy invariant (always relax the nearest unfinished vertex) is correct only when all weights are non-negative. On \texttt{roads.csv} the optimal route to~$B$ uses the cheaper detour $A\to C\to B$ (cost 3) instead of the direct edge $A\to B$ (cost 4).
  \item \textbf{Bellman-Ford} runs in $O(VE)$. It performs $|V|-1$ full relaxation passes and then one extra pass to check for negative cycles. It is slower but handles negative-weight edges and detects negative cycles.
  \item \textbf{When to prefer each.} Use Dijkstra for large sparse graphs without negative weights (e.g., road-navigation with positive distances). Use Bellman-Ford when negative weights appear (e.g., financial graphs with gains and losses) or when negative-cycle detection is required.
\end{itemize}

% ─────────────────────────────────────────────────────────────────────────────
\subsection{B4 --- Floyd-Warshall All-Pairs Shortest Paths}

\begin{lstlisting}[style=python, caption={Floyd-Warshall with path reconstruction (src/b4\_floyd\_warshall.py).}]
def floyd_warshall(vertices, edges):
    n   = len(vertices)
    idx = {v: i for i, v in enumerate(vertices)}

    dist = [[INF]  * n for _ in range(n)]
    nxt  = [[None] * n for _ in range(n)]

    for i in range(n):
        dist[i][i] = 0

    for u, v, w in edges:
        i, j = idx[u], idx[v]
        if w < dist[i][j]:
            dist[i][j] = w
            nxt[i][j]  = j     # next-hop index

    for k in range(n):
        for i in range(n):
            for j in range(n):
                via = dist[i][k] + dist[k][j]
                if via < dist[i][j]:
                    dist[i][j] = via
                    nxt[i][j]  = nxt[i][k]   # inherit next hop

    has_neg_cycle = any(dist[i][i] < 0 for i in range(n))
    return dist, nxt, has_neg_cycle, idx


def reconstruct_path(vertices, nxt, idx, src, dst):
    i, j = idx[src], idx[dst]
    if nxt[i][j] is None:
        return None
    path = [src]
    while i != j:
        i = nxt[i][j]
        path.append(vertices[i])
    return path
\end{lstlisting}

\paragraph{Output.}
\begin{lstlisting}[style=output]
Floyd-Warshall All-Pairs Shortest Paths
==================================================
Negative cycle detected: False

Distance matrix (INF = no path):
      A       B       C       D
A     0       3       2       8
B   INF       0     INF       5
C   INF       1       0       6
D   INF     INF     INF       0

Path reconstruction (all reachable pairs):
------------------------------------------------
A -> B: dist=3, path=A -> C -> B
A -> C: dist=2, path=A -> C
A -> D: dist=8, path=A -> C -> B -> D
B -> A: unreachable
B -> C: unreachable
B -> D: dist=5, path=B -> D
C -> A: unreachable
C -> B: dist=1, path=C -> B
C -> D: dist=6, path=C -> B -> D
D -> A/B/C: unreachable
\end{lstlisting}

\paragraph{Path reconstruction.}
The \texttt{nxt[i][j]} table stores the index of the next vertex after $i$ on the shortest path to $j$. To reconstruct path $A\to D$: $\mathtt{nxt}[A][D]=C$, then $\mathtt{nxt}[C][D]=B$, then $\mathtt{nxt}[B][D]=D$ -- yielding $A\to C\to B\to D$ with cost 8.

\paragraph{Negative-cycle detection.}
If \texttt{dist[i][i] < 0} after the algorithm completes, vertex $i$ lies on a negative-weight cycle. No such cycle exists here.

\paragraph{Reflection: Floyd-Warshall vs.\ repeated Dijkstra.}
\begin{itemize}[noitemsep]
  \item Floyd-Warshall is $O(V^3)$, simple to implement, and directly produces the full $V\times V$ distance matrix. It is a good choice when $V$ is small (here $V=4$) or the graph is dense.
  \item Repeated Dijkstra (once per source) is $O(V\cdot(V+E)\log V)$. For large sparse non-negative-weight graphs, this is typically faster.
  \item Floyd-Warshall is \emph{not} a good choice when $V$ is large (e.g., $V=100{,}000$ would require $10^{15}$ operations) or when only a few source-to-destination pairs are needed.
\end{itemize}

% ═══════════════════════════════════════════════════════════════════════════════
\section{Part C --- Greedy Network Design}
% ═══════════════════════════════════════════════════════════════════════════════

Part C uses \texttt{data/infra.csv}: an undirected weighted graph with vertices $\{A,B,C,D\}$ and edges $A\text{-}B\ (7)$, $A\text{-}C\ (2)$, $B\text{-}C\ (5)$, $B\text{-}D\ (6)$, $C\text{-}D\ (4)$.

% ─────────────────────────────────────────────────────────────────────────────
\subsection{C1 --- MST via Prim's Algorithm}

\begin{lstlisting}[style=python, caption={Prim's MST (src/c1\_prim.py).}]
def prim(graph, vertices, start):
    in_mst = set()
    key    = {v: float('inf') for v in vertices}
    parent = {v: None         for v in vertices}
    key[start] = 0
    pq = [(0, start)]

    mst_edges  = []
    total_cost = 0

    while pq:
        cost, u = heapq.heappop(pq)
        if u in in_mst:       # stale heap entry
            continue
        in_mst.add(u)

        if parent[u] is not None:
            mst_edges.append((parent[u], u, cost))
            total_cost += cost

        for v, w in graph.get(u, []):
            if v not in in_mst and w < key[v]:
                key[v]    = w
                parent[v] = u
                heapq.heappush(pq, (w, v))

    return mst_edges, total_cost
\end{lstlisting}

\paragraph{Output.}
\begin{lstlisting}[style=output]
Prim's Minimum Spanning Tree
===================================
Starting vertex: A

Edge            Weight
-------------------------
A -- C          2
C -- D          4
C -- B          5
-------------------------
Total MST cost: 11

MST spans all 4 vertices (3 edges) -- valid spanning tree.
All MST edge weights are distinct -> the MST is unique.
\end{lstlisting}

\paragraph{MST uniqueness.}
When all edge weights are distinct, the MST is unique (provable by contradiction using the cut property). All five edges in \texttt{infra.csv} have distinct weights (2, 4, 5, 6, 7), so the MST $\{A\text{-}C,\, C\text{-}D,\, C\text{-}B\}$ with cost 11 is the only MST.

Had two edges tied for the same weight, multiple valid MSTs could exist; in that case we would report both and note they share the same total cost.

% ─────────────────────────────────────────────────────────────────────────────
\subsection{C2 --- Greedy Correctness Discussion}

\paragraph{Why MST is a greedy problem.}
MST algorithms are greedy in the sense that each step makes a locally optimal choice -- add the cheapest edge that does not violate the spanning-tree constraints -- and this local choice is provably globally optimal. The correctness rests on two structural properties:
\begin{itemize}[noitemsep]
  \item \textbf{Cut property.} For any cut $(S, V\setminus S)$ of the graph, the minimum-weight crossing edge belongs to some MST. Prim's algorithm always adds the cheapest edge that crosses the frontier between the in-MST and not-yet-in-MST vertices, respecting this property at every step.
  \item \textbf{Cycle property.} The maximum-weight edge in any cycle cannot belong to any MST (it can always be removed without disconnecting the tree).
\end{itemize}

\paragraph{Why na\"ive ``pick the globally cheapest remaining edge'' fails without cycle checking.}
Consider the graph $A\text{-}B\ (2)$, $B\text{-}C\ (3)$, $A\text{-}C\ (4)$. A valid MST contains edges $\{A\text{-}B,\, B\text{-}C\}$ (cost 5).

If we pick the cheapest remaining edge at each step \emph{without checking for cycles}:
\begin{enumerate}[noitemsep]
  \item Pick $A\text{-}B\ (2)$. Selected: $\{A\text{-}B\}$.
  \item Pick $B\text{-}C\ (3)$. Selected: $\{A\text{-}B, B\text{-}C\}$.
  \item Pick $A\text{-}C\ (4)$. \textbf{Without cycle checking}, we add this edge too.
\end{enumerate}
Result: $\{A\text{-}B, B\text{-}C, A\text{-}C\}$, which is a \emph{cycle} -- no longer a tree. Cycle checking (as done by Kruskal's algorithm using a Union-Find structure) is essential to ensure the result remains a tree.

\paragraph{Counterexample (concrete).}
Use the same three-vertex graph above. The na\"ive greedy without cycle detection selects all three edges ($A\text{-}B$, $B\text{-}C$, $A\text{-}C$) at costs 2, 3, 4 = total 9. A valid spanning tree needs only two edges; the MST has cost 5. The na\"ive result is both more expensive and structurally invalid (it contains a cycle).

% ═══════════════════════════════════════════════════════════════════════════════
\section{Part D --- Dynamic Programming Scheduling}
% ═══════════════════════════════════════════════════════════════════════════════

Part D uses \texttt{data/jobs.csv}:

\begin{center}
\begin{tabular}{lrrr}
\toprule
Job & Start & Finish & Reward \\
\midrule
J1 & 1 & 4 & 5 \\
J2 & 3 & 5 & 2 \\
J3 & 0 & 6 & 8 \\
J4 & 4 & 7 & 6 \\
J5 & 3 & 8 & 7 \\
\bottomrule
\end{tabular}
\end{center}

% ─────────────────────────────────────────────────────────────────────────────
\subsection{D1 --- Greedy Interval Scheduling (Maximise Job Count)}

\begin{lstlisting}[style=python, caption={Greedy earliest-finish-time scheduling (src/d1\_greedy\_scheduling.py).}]
def greedy_interval_scheduling(jobs):
    sorted_jobs  = sorted(jobs, key=lambda j: j['finish'])
    selected     = []
    last_finish  = -1    # sentinel

    for job in sorted_jobs:
        if job['start'] >= last_finish:
            selected.append(job)
            last_finish = job['finish']

    return selected
\end{lstlisting}

\paragraph{Output.}
\begin{lstlisting}[style=output]
All jobs sorted by finish time:
Job      Start    Finish   Reward
------------------------------------
J1       1        4        5
J2       3        5        2
J3       0        6        8
J4       4        7        6
J5       3        8        7

Selected jobs (2 non-overlapping):
Job      Start    Finish   Reward
------------------------------------
J1       1        4        5
J4       4        7        6

Jobs selected : 2
Total reward  : 11
\end{lstlisting}

The greedy algorithm selects J1 (finishes earliest at time 4) then J4 (starts exactly at 4, compatible), yielding the maximum possible count of 2 non-overlapping jobs. J2, J3, and J5 all overlap with J1.

% ─────────────────────────────────────────────────────────────────────────────
\subsection{D2 --- Weighted Interval Scheduling (DP)}

\paragraph{Recurrence.}
Sort jobs by finish time and index them $1 \ldots n$. Let $p(j)$ be the index of the latest job $i < j$ with $\mathrm{finish}(i) \le \mathrm{start}(j)$ (0 if none). Define:
\[
\mathrm{OPT}(0) = 0, \qquad
\mathrm{OPT}(j) = \max\!\bigl(\,\mathrm{reward}_j + \mathrm{OPT}(p(j)),\;\mathrm{OPT}(j-1)\bigr).
\]
The first argument includes job $j$; the second excludes it.

\begin{lstlisting}[style=python, caption={Predecessor computation and DP (src/d2\_weighted\_dp.py).}]
def compute_predecessors(jobs):
    n = len(jobs)
    p = [-1] * n
    for j in range(n):
        for i in range(j - 1, -1, -1):
            if jobs[i]['finish'] <= jobs[j]['start']:
                p[j] = i
                break
    return p


def weighted_interval_scheduling(jobs):
    jobs_sorted = sorted(jobs, key=lambda j: j['finish'])
    n  = len(jobs_sorted)
    p  = compute_predecessors(jobs_sorted)

    dp = [0] * (n + 1)
    for j in range(1, n + 1):
        job         = jobs_sorted[j - 1]
        pred_dp_idx = p[j - 1] + 1          # shift 0-indexed p to 1-indexed dp
        dp[j] = max(job['reward'] + dp[pred_dp_idx], dp[j - 1])

    # Reconstruction: walk backwards
    selected = []
    j = n
    while j >= 1:
        job         = jobs_sorted[j - 1]
        pred_dp_idx = p[j - 1] + 1
        if job['reward'] + dp[pred_dp_idx] >= dp[j - 1]:
            selected.append(job)
            j = pred_dp_idx
        else:
            j -= 1
    selected.reverse()
    return dp[n], selected, dp, jobs_sorted, p
\end{lstlisting}

\paragraph{Output.}
\begin{lstlisting}[style=output]
Jobs sorted by finish time (p uses 1-based indexing, 0 = none):
j     Job      Start    Finish   Reward   p(j)
------------------------------------------------
1     J1       1        4        5        0
2     J2       3        5        2        0
3     J3       0        6        8        0
4     J4       4        7        6        1
5     J5       3        8        7        0

Recurrence:
OPT(0) = 0
OPT(j) = max( reward_j + OPT(p(j)),  OPT(j-1) )

DP table:
j      OPT(j)
--------------
0      0
1      5
2      5
3      8
4      11
5      11

Maximum total reward: 11

Selected jobs:
Job      Start    Finish   Reward
------------------------------------
J1       1        4        5
J4       4        7        6

Total reward: 11
\end{lstlisting}

\paragraph{DP table walkthrough.}
\begin{itemize}[noitemsep]
  \item $\mathrm{OPT}(1) = \max(5+0,\, 0) = 5$ \quad (include J1).
  \item $\mathrm{OPT}(2) = \max(2+0,\, 5) = 5$ \quad (exclude J2; its reward alone is too small).
  \item $\mathrm{OPT}(3) = \max(8+0,\, 5) = 8$ \quad (include J3 alone gives 8 $>$ 5).
  \item $\mathrm{OPT}(4) = \max(6+\mathrm{OPT}(1),\, \mathrm{OPT}(3)) = \max(11,8) = 11$ \quad (J4 is compatible with J1).
  \item $\mathrm{OPT}(5) = \max(7+0,\, 11) = 11$ \quad (exclude J5).
\end{itemize}

\paragraph{Reconstruction.}
Starting from $j=5$: $7 < 11$, exclude J5, $j=4$. At $j=4$: $11 \ge 8$, include J4, jump to $j=p(4)=1$. At $j=1$: $5 \ge 0$, include J1, jump to $j=0$. Stop. Selected: J1, J4, total reward 11.

% ─────────────────────────────────────────────────────────────────────────────
\subsection{D3 --- Greedy Failure Example for Weighted Scheduling}

\paragraph{Counterexample.}
Consider two jobs:
\begin{center}
\begin{tabular}{lrrr}
\toprule
Job & Start & Finish & Reward \\ \midrule
$A$ & 1 & 2 & 1 \\
$B$ & 1 & 10 & 100 \\ \bottomrule
\end{tabular}
\end{center}

Both jobs start at time 1 and are mutually incompatible (overlapping).

\begin{itemize}
  \item \textbf{Greedy (earliest finish, D1 strategy).} Sorts: $A$ finishes at 2, $B$ at 10. Selects $A$ (finishes earlier); then $B$ starts at 1 $< 2$, so skip. \textbf{Result:} $\{A\}$, count = 1, reward = 1.
  \item \textbf{DP (weighted, D2 strategy).} $\mathrm{OPT}(1)=\max(1,0)=1$ (take $A$). $\mathrm{OPT}(2)=\max(100+0,\,1)=100$ (take $B$). \textbf{Result:} $\{B\}$, reward = 100.
\end{itemize}

\paragraph{Explanation of the difference.}
The two objectives are fundamentally different:
\begin{itemize}[noitemsep]
  \item \textbf{Maximise job count.} Greedy by earliest finish is optimal: choosing the job that finishes soonest leaves the maximum remaining time for future jobs.
  \item \textbf{Maximise total reward.} Greedy by earliest finish can be arbitrarily bad: it may pick a low-reward job that finishes early, blocking a high-reward job. In the example above, greedy earns reward 1 while the optimal earns 100 -- a factor-of-100 loss.
\end{itemize}

Dynamic programming is required for the weighted version because the optimal solution may skip short, cheap jobs in favour of longer, more rewarding ones -- a trade-off that greedy cannot capture.

\end{document}
